% !TeX program = xelatex
\documentclass[12pt]{article}
\usepackage[spanish, es-nodecimaldot]{babel}

\usepackage{fontspec}
\setmainfont{Times New Roman}

\usepackage[
letterpaper,
top=2.4cm,
bottom=2.4cm,
left=2.4cm,
right=2.4cm
]{geometry}


\usepackage{setspace}
% Renglon cerrado
\singlespacing

\usepackage[hidelinks]{hyperref}

\usepackage{enumitem}

\setlength{\parindent}{0pt}

\setlength{\parskip}{6pt}

\author{Jenny Sofía Morales}

\begin{document}
	
	\begin{center}
		
		{\bfseries\fontsize{12}{14}\selectfont
			LA ENCUESTA: PREGUNTAR CON PROPOSITO
		}
		
		\vspace{0.4cm}
		
		{\itshape J. S. Morales López}
		
		\vspace{0.2cm}
		
		{\itshape 7690-08-6790 Universidad Mariano Gálvez}
		
		{\itshape Seminario de tecnologías de información}
		
		{\itshape
			\href{mailto:jmorales115@miumg.edu.gt}
			{jmorales115@miumg.edu.gt}
		}
		
	\end{center}
	
	\vspace{0.5cm}
	\section*{Resumen}
	
	El presente artículo analiza los elementos necesarios para diseñar y
	aplicar una encuesta de manera efectiva. Su objetivo es explicar cómo la
	definición del propósito, la selección de la población y la muestra, la
	redacción de preguntas, la elección del medio de aplicación y el
	análisis de los datos influyen en la calidad de la información recopilada.
	Para desarrollar una encuesta se investigaron conceptos relacionados con la
	planificación del cuestionario, el lenguaje utilizado, la duración de la
	encuesta, la privacidad de los participantes y el uso responsable de los
	resultados. A partir del análisis realizado, se identificó que una encuesta
	es efectiva cuando permite obtener información útil para responder al
	objetivo planteado, y es óptima cuando alcanza ese resultado mediante
	preguntas previamente planificadas y un uso adecuado del tiempo y los recursos
	disponibles. También se identificó que una mala selección de participantes,
	preguntas ambiguas o un cuestionario demasiado largo pueden reducir
	la confiabilidad de las respuestas. Se puede decir que la encuesta no debe
	elaborarse de manera improvisada, sino debe llevar un proceso organizado que
	permita convertir las respuestas de los participantes en información
	clara, confiable y útil para la toma de decisiones.
	
	\noindent
	\textbf{Palabras clave:} encuesta, cuestionario, muestra, efectiva, óptima,	resultados.
	
	
	\section*{Desarrollo del tema}
	
	La encuesta es uno de los recursos más utilizados dentro de una
	investigación, debido a que permite recopilar información directamente
	de un grupo de personas acerca de un tema específico. Por medio de
	preguntas planificadas con anterioridad, el investigador puede conocer
	opiniones, experiencias, necesidades, comportamientos o niveles de
	satisfacción.
	
	Sin embargo, aplicar una encuesta no consiste únicamente en formular
	preguntas y esperar respuestas. Para que los resultados sean útiles,
	debe existir una planificación adecuada que considere el objetivo del
	estudio, las características del público, el lenguaje empleado, el tipo
	de preguntas y la forma en que se analizarán los datos.
	
	Según Linares Fontela (2003), la elaboración de una encuesta requiere
	un proceso organizado que incluye la selección del tipo de encuesta,
	la definición de la población y la muestra, el diseño del cuestionario,
	su aplicación y el análisis de los resultados.
	
	Una encuesta puede considerarse efectiva cuando permite obtener la
	información necesaria para cumplir el objetivo de la investigación. Por
	otra parte, puede considerarse óptima cuando obtiene esa información
	mediante un cuestionario claro, breve, ordenado y adecuado para las
	personas participantes o publico objetivo.
	
	El primer paso consiste en definir el objetivo de la encuesta. Antes de
	redactar las preguntas, se debe establecer qué información se desea
	obtener y para qué será utilizada. Cuando el objetivo es demasiado
	amplio o ambiguo, pueden incluirse preguntas innecesarias que dificulten el
	análisis de los resultados.
	
	También es importante seleccionar adecuadamente la población y la muestra. La población corresponde al conjunto de personas relacionadas con el tema de investigación, mientras que la muestra representa al grupo específico que responderá la encuesta.
	
	Las preguntas deben redactarse con un lenguaje sencillo, claro y comprensible. Se deben evitar palabras demasiado técnicas, preguntas	muy largas o expresiones que puedan influir en las respuestas de los	participantes.
	
	La preguntas pueden ser abiertas, cerradas, de opción multiple o de escala.	Las preguntas abiertas permiten que los participantes respondan con sus propias palabras y su opinion detallada.	Las preguntas cerradas dan un tono mas controlado del tema con respuestas predefinidas para facilitar la obtención de los resultados.	Por otra parte las preguntas a escala permiten medir por ejemplo el nivel de satisfacción, o acuerdos. 	La elección del tipo de preguntas depende de el objetivo de la encuesta y los datos que necesita obtener el investigador.
	
	El medio para realizar una encuesta también debe ser cuidadosamente considerado, ya que existe distintas modalidades. En las encuestas telefónicas debe considerarse el tono de voz del encuestador y la claridad para que sea efectiva.	En la encuesta física o escrita debe considerarse el lugar donde se realiza y los materiales para realizarla. Las encuestas digitales debe considerarse el publico objetivo porque algunas personas no cuentan con el acceso al internet o no estar familiarizado con las herramientas tecnologícas.
	
	Finalmente, los resultados deben organizarse, analizarse e
	interpretarse. No es suficiente presentar porcentajes o gráficas; el
	investigador debe explicar qué representan los datos y cómo se
	relacionan con el objetivo planteado.
	
	\section*{Observaciones y comentarios}
	
	Durante la elaboración de una encuesta se debe prestar especial atención a la duración del cuestionario. Una encuesta demasiado larga puede ser aburrida, y provocar desinteres. También es importante respetar la privacidad de los participantes. Los datos personales deben solicitarse únicamente cuando sean necesarios e indicarle a los participantes el porque y para qué es la encuesta.
	
	\section*{Conclusiones}
	
	\begin{enumerate}
		
		\item Una encuesta efectiva requiere un objetivo claro, una muestra adecuada y preguntas claras que permitan obtener información relacionada con el propósito de la investigación.
		
		\item Una encuesta puede considerarse óptima cuando alcanza el objetivo planteado mediante un cuestionario no improvisado y un uso eficiente del tiempo y los recursos disponibles.
		
		\item La planificación, la neutralidad de las preguntas y el análisis
		responsable de los resultados permiten transformar las respuestas en
		información confiable para la toma de decisiones.
		
		\item La identificación clara del publico objetivo permite redactar preguntas adecuadas con las caracteristicas necesarias, lo que contribuye a la obtener información confiable.
		
	\end{enumerate}
	
	\section*{Bibliografía}
	
	\begin{itemize}[leftmargin=0cm, label={}]
		
		\item Hernández Sampieri, R., Fernández Collado, C. y Baptista
		Lucio, M. P. (2014). \textit{Metodología de la investigación}.
		McGraw-Hill.
		
		\item Martínez, F. (2002). El cuestionario: un instrumento para la
		investigación en las ciencias sociales.
		\textit{Revista de Investigación Educativa}.
		
		\item Organización de las Naciones Unidas para la Educación, la
		Ciencia y la Cultura. (2021). \textit{Métodos para la recopilación
			y el análisis de información}.
		\item Linares Fontela, J. (2003).
		\textit{Guía para diseñar encuestas}.
		En B. Branch y J. Klaehn (Eds.),
		\textit{El logro del equilibrio en las microfinanzas:
			Guía práctica para movilizar los ahorros}
		(pp. 327--336). Pact Publications.
		\url{https://www.woccu.org/documents/Tool10(sp)}
		
	\end{itemize}
	
\end{document}